\section{Architettura Logica}
L'architettura logica del progetto SmartOrder adotta un approccio esagonale.\\
L'obiettivo è quello di mantenere separata la logica di business dalle componenti esterne, come l'interfaccia utente e il sistema di persistenza dei dati.\\
La comunicazione tra la parte di presentazione e la logica applicativa avviene tramite interfacce, che permettono ai controller (ad esempio quelli esposti tramite API HTTP) di invocare le funzionalità del sistema senza dipendere direttamente dalla loro implementazione.\\
Allo stesso modo, l'accesso al database e ad altri servizi esterni è gestito tramite componenti dedicati, che permettono alla logica di business di rimanere indipendente dalle tecnologie utilizzate.

\subsection{Pattern di architettura esagonale}
Il pattern esagonale è stato scelto per organizzare il sistema in modo chiaro e modulare, separando le responsabilità tra le diverse componenti.\\
Questo approccio rende il sistema più facile da mantenere, testare ed estendere nel tempo.\\
I principali elementi dell'architettura esagonale sono:\\
\begin{itemize}
	\item \textbf{Core domain}: contiene la logica di business dell'applicazione ed è indipendente dalle tecnologie esterne.
	\item \textbf{Porte}: rappresentano i punti di contatto tra il dominio e l'esterno, e definiscono le operazioni disponibili. Si distinguono in:
	\begin{itemize}
		\item \textbf{Inbound ports}: interfacce che espongono le funzionalità del dominio verso l'esterno
		\item \textbf{Outbound ports}: interfacce che definiscono come il dominio può utilizzare servizi esterni
	\end{itemize}
	\item \textbf{Adapters}: sono implementazioni concrete delle porte che collegano il dominio alle tecnologie utilizzate
	\begin{itemize}
		\item \textbf{Inbound adapters}: adattatori che convertono le richieste esterne nel formato atteso del dominio
		\item \textbf{Outbound adapters}: adattori che implementano le interfacce di uscita collegandole a tecnologie specifiche
	\end{itemize}
\end{itemize}

\section{Servizi principali}
\begin{itemize}
	\item Autenticazione e registrazione
\begin{figure}[H]
	\centering
	\includegraphics[width=0.8\textwidth]{SpecificaTecnica/auth}
	\caption{Registrazione e autenticazione}
\end{figure}
	\begin{itemize}
		\item \textbf{Core domain}: Contiene la logica dedicata alla registrazione di un utente nuovo oppure all'autenticazione di un utente già esistente. Permette di elaborare i dati ricevuti tramite l'inserimento dei dati dell'utente.
		\item \textbf{Inbound ports}: Non sono presenti inbound ports esplicite. Il layer applicativo è esposto direttamente tramite \texttt{UserService}, che viene istanziato e invocato dal controller FastAPI. Questa scelta è giustificata dalla presenza di un unico driver (l'interfaccia HTTP), che rende superflua l'introduzione di un'interfaccia dedicata sul lato entrante.
		\item \textbf{Outbound ports}: Definiscono i contratti verso l'infrastruttura esterna. \texttt{IUserRepoPort} espone le operazioni di persistenza sugli utenti, mentre \texttt{IEmailValidationPort} espone la verifica del dominio email tramite DNS. Entrambe le interfacce sono definite nel core domain e implementate da adapter dedicati.
		\item \textbf{Adapters}: Implementa adattatori per l'inserimento nel DB
	\end{itemize}
	\item Chatbot e gestione ordine
	\begin{itemize}
		\item \textbf{Core domain}:
		\item \textbf{Inbound ports}:
		\item \textbf{Outbound ports}:
		\item \textbf{Adapters}:
	\end{itemize}
	\item Visualizzazione storico ordini
	\begin{itemize}
		\item \textbf{Core domain}:
		\item \textbf{Inbound ports}:
		\item \textbf{Outbound ports}:
		\item \textbf{Adapters}:
	\end{itemize}
\end{itemize}